% Fragebogen der Anwenderbefragung (Abschnitt 7.3, Anhang).
% Eigenständiges Dokument: wird ausgedruckt in der Fertigung verteilt (ein Blatt, beidseitig) und als PDF
% in den Anhang eingebunden (\includepdf[pages=-]{Abbildungen/Fragebogen.pdf}, Paket pdfpages).
% Übersetzen (aus diesem Ordner):
%   latexmk -pdf -outdir=../build/Fragebogen Fragebogen.tex
%   cp ../build/Fragebogen/Fragebogen.pdf .
%
% Teil B ist die System Usability Scale (SUS) nach Brooke 1996 in der gebräuchlichen deutschen Fassung.
% Wortlaut und Reihenfolge nicht ändern, sonst ist der SUS-Wert nicht mehr vergleichbar.
% Auswertung je Bogen: ungerade Aussagen (Wert - 1), gerade Aussagen (5 - Wert), Summe mal 2,5 -> 0 bis 100.
\documentclass[11pt,a4paper]{article}
\usepackage[ngerman]{babel}
\usepackage[utf8]{inputenc}
\usepackage[T1]{fontenc}
\usepackage{helvet}
\renewcommand{\familydefault}{\sfdefault}
\usepackage[a4paper,top=1.4cm,bottom=1.3cm,left=1.7cm,right=1.7cm]{geometry}
\usepackage[table]{xcolor}
\usepackage{array,tabularx}
\usepackage{amssymb,graphicx}
\usepackage[babel,german=quotes]{csquotes}

\pagestyle{empty}
\setlength{\parindent}{0pt}
\setlength{\parskip}{0.35em}

\definecolor{zeile}{gray}{0.92}
% Ankreuzkästchen
\newcommand{\kasten}{\raisebox{-0.2ex}{\scalebox{1.5}{$\square$}}}
% Schreiblinie über die ganze Breite
\newcommand{\linie}{\par\vspace{0.6em}\noindent\rule{\linewidth}{0.4pt}}
\newcommand{\teil}[1]{\par\vspace{0.5em}{\large\bfseries #1}\par}

% Spalten für die Antwortskala
\newcolumntype{K}{>{\centering\arraybackslash}m{0.85cm}}
\newcolumntype{N}{>{\raggedleft\arraybackslash}m{0.75cm}}
\newcolumntype{A}{>{\raggedright\arraybackslash}X}

% Kopf der Skala (eine Tabellenzeile)
\newcommand{\eins}{\makebox[0pt]{\shortstack{\bfseries 1\\[-0.1em]\scriptsize gar nicht}}}
\newcommand{\fuenf}{\makebox[0pt]{\shortstack{\bfseries 5\\[-0.1em]\scriptsize voll}}}
\newcommand{\zeile}[2]{\bfseries #1 & #2 & \kasten & \kasten & \kasten & \kasten & \kasten \\}
\newcommand{\zeileX}[2]{\bfseries #1 & #2 & \kasten & \kasten & \kasten & \kasten & \kasten & \kasten \\}

\begin{document}

{\LARGE\bfseries Befragung zur Produktionsverwaltung}

\vspace{0.3em}
Liebe Kolleginnen und Kollegen,

im Rahmen meiner Bachelorarbeit an der HTWK Leipzig untersuche ich, wie gut sich die
Produktionsverwaltung im Arbeitsalltag bewährt. Ihre Erfahrungen helfen, die Anwendung zu
bewerten und weiter zu verbessern.

Die Teilnahme ist \textbf{freiwillig und anonym}. Bitte schreiben Sie Ihren Namen nicht auf den
Bogen. Die Antworten werden nur zusammengefasst ausgewertet, Rückschlüsse auf einzelne Personen
sind nicht möglich. Das Ausfüllen dauert etwa \textbf{zehn Minuten}.

Bitte geben Sie den ausgefüllten Bogen bis \textbf{Freitag, 04.12.2026} ab bei:
\rule{5.2cm}{0.4pt}

Vielen Dank für Ihre Unterstützung!\\
Arvid Graatz

%--------------------------------------------------------------------------
\teil{Teil A: Ihre Tätigkeit}

\textbf{A1}\enspace In welchen Bereichen arbeiten Sie mit der Produktionsverwaltung?
(Mehrfachnennung möglich)

\begin{tabular}{@{\hspace{1.8em}}l@{\hspace{2em}}l@{\hspace{2em}}l}
	\kasten\ Bestückung (SMD/THT) & \kasten\ Prüfung / optische Kontrolle & \kasten\ Endkontrolle \\[0.35em]
	\kasten\ Lager und Versand    & \kasten\ Modulmontage                 & \kasten\ Planung / Leitung \\[0.35em]
	\multicolumn{3}{@{\hspace{1.8em}}l}{\kasten\ Sonstiges: \rule{6cm}{0.4pt}}
\end{tabular}

\textbf{A2}\enspace Wie oft nutzen Sie die Produktionsverwaltung?

\begin{tabular}{@{\hspace{1.8em}}l@{\hspace{2em}}l@{\hspace{2em}}l@{\hspace{2em}}l}
	\kasten\ täglich & \kasten\ mehrmals pro Woche & \kasten\ etwa einmal pro Woche & \kasten\ seltener
\end{tabular}

\textbf{A3}\enspace Haben Sie vor der Produktionsverwaltung mit Papier-Laufkarten oder Excel-Listen gearbeitet?

\begin{tabular}{@{\hspace{1.8em}}l@{\hspace{2em}}l}
	\kasten\ ja & \kasten\ nein
\end{tabular}

%--------------------------------------------------------------------------
\teil{Teil B: Bedienung allgemein}

Bitte kreuzen Sie bei jeder Aussage spontan an, wie sehr Sie zustimmen (1 = stimme gar nicht zu,
5 = stimme voll zu). Mit \enquote{System} ist die Produktionsverwaltung gemeint. Bitte beantworten
Sie \textbf{alle} Aussagen. Wenn Sie unsicher sind, kreuzen Sie die Mitte (3) an.

{\renewcommand{\arraystretch}{1.55}
\rowcolors{2}{zeile}{white}
\begin{tabularx}{\linewidth}{@{}N@{\hspace{0.6em}}AKKKKK@{}}
	& & \eins & \bfseries 2 & \bfseries 3 & \bfseries 4 & \fuenf \\
	\zeile{B1}{Ich denke, dass ich das System gerne häufig benutzen würde.}
	\zeile{B2}{Ich fand das System unnötig komplex.}
	\zeile{B3}{Ich fand das System einfach zu benutzen.}
	\zeile{B4}{Ich glaube, ich würde die Hilfe einer technisch versierten Person benötigen, um das System benutzen zu können.}
	\zeile{B5}{Ich fand, die verschiedenen Funktionen in diesem System waren gut integriert.}
	\zeile{B6}{Ich denke, das System enthielt zu viele Inkonsistenzen.}
	\zeile{B7}{Ich kann mir vorstellen, dass die meisten Menschen den Umgang mit diesem System sehr schnell lernen.}
	\zeile{B8}{Ich fand das System sehr umständlich zu nutzen.}
	\zeile{B9}{Ich fühlte mich bei der Benutzung des Systems sehr sicher.}
	\zeile{B10}{Ich musste eine Menge lernen, bevor ich anfangen konnte, mit dem System zu arbeiten.}
\end{tabularx}}

\newpage
%--------------------------------------------------------------------------
\teil{Teil C: Einzelne Funktionen und Vergleich mit früher}

Kreuzen Sie \enquote{kann ich nicht beurteilen} an, wenn Sie eine Funktion nicht nutzen oder
(bei C11 bis C13) früher nicht mit Papier-Laufkarten oder Excel gearbeitet haben.

{\renewcommand{\arraystretch}{1.3}
\rowcolors{2}{zeile}{white}
\begin{tabularx}{\linewidth}{@{}N@{\hspace{0.6em}}AKKKKK>{\centering\arraybackslash}m{1.9cm}@{}}
	& & \eins & \bfseries 2 & \bfseries 3 & \bfseries 4 & \fuenf & \shortstack{\scriptsize kann ich nicht\\[-0.1em]\scriptsize beurteilen} \\
	\zeileX{C1}{Die Erfassung von Produktionsschritten mit dem Scanner geht schnell.}
	\zeileX{C2}{Die Erfassung über einen Seriennummernbereich ist einfach.}
	\zeileX{C3}{Wenn eine Erfassung abgelehnt wird (z.\,B. weil ein Schritt fehlt), verstehe ich den Grund.}
	\zeileX{C4}{Ich finde schnell heraus, in welchem Schritt sich eine Platine gerade befindet.}
	\zeileX{C5}{Ich kann gut nachvollziehen, was mit einer bestimmten Platine passiert ist.}
	\zeileX{C6}{Wenn ich mich vertan habe, kann ich den Fehler über eine Korrekturanfrage gut melden.}
	\zeileX{C7}{Die Seiten laden schnell genug für meine Arbeit.}
	\zeileX{C8}{Platinen in Kisten oder auf Paletten packen, versenden und annehmen ist verständlich.}
	\zeileX{C9}{Platinen per Scan in Module einbauen ist einfach.}
	\zeileX{C10}{Die Benachrichtigungen (z.\,B. zu Lieferterminen) sind nützlich.}
	\zeileX{C11}{Im Vergleich zu Papier-Laufkarten und Excel spare ich bei der Erfassung Zeit.}
	\zeileX{C12}{Im Vergleich zu früher passieren bei der Erfassung weniger Fehler.}
	\zeileX{C13}{Informationen zu einer Platine finde ich schneller als früher.}
\end{tabularx}}

%--------------------------------------------------------------------------
\teil{Teil D: Ihre Meinung}

\textbf{D1}\enspace Was gefällt Ihnen an der Produktionsverwaltung am besten?
\linie\linie\linie

\textbf{D2}\enspace Was stört Sie, oder was fehlt Ihnen?
\linie\linie\linie

\textbf{D3}\enspace Mussten Sie schon einmal auf Papier oder Excel ausweichen? Wenn ja, in welcher Situation?
\linie\linie

\vfill
\begin{center}
	\bfseries Vielen Dank für Ihre Zeit!
\end{center}

\end{document}
